% concatenated from JS0822.tex
\documentclass[11pt]{amsart}
%\usepackage[scrtime]{prelim2e}
\usepackage{amsmath,amssymb,amsthm,latexsym,cite,cancel}
\usepackage[small]{caption}
\usepackage{graphicx,wasysym,overpic,tikz,color}
\usepackage{subfigure}
\usepackage{cite}
\usepackage[colorlinks=true,urlcolor=blue,
citecolor=red,linkcolor=blue,linktocpage,pdfpagelabels,
bookmarksnumbered,bookmarksopen]{hyperref}
\usepackage[english]{babel}
\usepackage{units}
\usepackage{enumitem}
\usepackage[left=2.1cm,right=2.1cm,top=2.71cm,bottom=2.71cm]{geometry}
\usepackage[hyperpageref]{backref}
\usepackage{float,bm}

%\usepackage{refcheck}
%\usepackage{showkeys}
\usepackage[colorinlistoftodos]{todonotes}
\makeatletter
\def\@seccntformat#1{%
	\protect\textup{\protect\@secnumfont
		\csname the#1\endcsname\protect\enspace}}
\makeatother
%
%\makeatletter
%\providecommand\@dotsep{5}
%\def\listtodoname{List of Todos}
%\def\listoftodos{\@starttoc{tdo}\listtodoname}
%\makeatother
%%%

%\usepackage[latin1]{inputenc}
%\usepackage{amsfonts}
%\usepackage{amscd}
%\usepackage{latexsym}
%\usepackage{mathrsfs}
%\usepackage[notref,notcite]{showkeys}
% \font\rmg=cmr10
%\font\bfg=cmbx10
%\renewcommand{\arraystretch}{1.5}
%\textwidth172mm \textheight22cm \hoffset-24mm \voffset-20mm
%\font\sans=cmss10
%\textwidth=13.5cm
%\baselineskip=16pt
%\makeatletter \@addtoreset{equation}{section} \makeatother
%\renewcommand{\theequation}{\arabic{section}.\arabic{equation}}
%\setlength{\parindent}{1em}

\newtheorem{theorem}{Theorem}[section]

\theoremstyle{definition}
\newtheorem{definition}[theorem]{Definition}

\theoremstyle{plain}
\newtheorem{proposition}[theorem]{Proposition}
\newtheorem{lemma}[theorem]{Lemma}
\newtheorem{remark}[theorem]{Remark}
\newtheorem{example}[theorem]{Example}
\newtheorem{corollary}[theorem]{Corollary}
\newtheorem{problem}[theorem]{Open Problem}
\newtheorem*{oldbrezisproblem}{Open Problem 1.1}
\newtheorem*{brezisproblem}{Brezis' Open Problem 2.1 \cite{Brezis2023}}
%\usepackage{CJK,CJKnumb,CJKulem,times,dsfont,ifthen,mathrsfs,latexsym,amsfonts,color}
%\usepackage{amsmath,amsthm,makeidx,fontenc,amssymb,bm,graphicx,psfrag,listings,curves,extarrows}

\newcommand{\abs}[1]{\lvert#1\rvert}
\newcommand{\norm}[1]{\lVert#1\rVert}
\newcommand{\red}[1]{\textcolor{red}{#1}}
\newcommand{\blue}[1]{\textcolor{blue}{#1}}
\newcommand{\D}{\mathcal{D}}
\newcommand{\G}{\mathcal{G}}
\newcommand{\K}{\mathcal{K}}
\newcommand{\h}{\mathcal{H}}
\newcommand{\E}{\mathrm{E}}
\newcommand{\V}{\mathrm{V}}
\newcommand{\R}{\mathbb{R}}
\newcommand{\C}{\mathbb{C}}
\newcommand{\ru}{\mathcal{u}}
\newcommand{\rv}{\mathcal{v}}

\tikzstyle{nodo}=[circle,draw,fill,inner sep=0pt,minimum size=%
1.5mm]

\numberwithin{equation}{section}
%\newtheorem{theorem}{Theorem}[section]
%\newtheorem{proposition}[theorem]{Proposition}
%\newtheorem{lemma}[theorem]{Lemma}
%\newtheorem{remark}[theorem]{Remark}
%\newtheorem{example}[theorem]{Example}
%\newtheorem{corollary}[theorem]{Corollary}
%\newtheorem{definition}[theorem]{Definition}
%\theoremstyle{definition}

\title[A Partial Answer to Brezis' Open Problem 2.1]
{A partial answer to Brezis' Open Problem 2.1}

\author[C. Ji]{Chao Ji}

\address[C. Ji]{\newline\indent
	School of Mathematics
	\newline\indent
	East China University of Science and Technology
	\newline\indent
	Shanghai 200237, PR China }
\email{\href{mailto:jichao@ecust.edu.cn}{jichao@ecust.edu.cn}}

\author[K. Sheng]{Kai Sheng}

\address[K. Sheng]{\newline\indent
	School of Mathematics
	\newline\indent
	East China University of Science and Technology
	\newline\indent
	Shanghai 200237, PR China}

\email{\href{mailto:shengkai2001@outlook.com}{shengkai2001@outlook.com}}
\subjclass[2020]{35Q56, 35A02, 35J57}
\date{\today}
\keywords{Ginzburg--Landau equation, uniqueness, convexity threshold, implicit function theorem}

\begin{document}
\begin{abstract}
In this paper, we give a partial answer to Brezis' Open Problem~2.1, which concerns the uniqueness of solutions to the Ginzburg--Landau equation in the unit disc with the degree-one boundary condition. Let $\lambda_1$ be the first Dirichlet eigenvalue of $-\Delta$ in the unit disc and set $\varepsilon_*:=\lambda_1^{-1/2}$. We prove that there exists $\delta>0$ such that the radial solution is the unique weak solution for every $\varepsilon\in(\varepsilon_*-\delta,\infty)$. More generally, we establish the analogous uniqueness result for the Ginzburg--Landau system on bounded connected $C^{1,1}$ domains in $\mathbb R^N$, $2\leq N\leq4$, with nontrivial boundary data. In particular, uniqueness persists slightly below the convexity threshold $\varepsilon_*$, where the strict convexity argument is no longer available. The proof combines strict convexity for $\varepsilon\geq\varepsilon_*$ with a compactness argument, nondegeneracy of the solution at $\varepsilon=\varepsilon_*$ and the implicit function theorem. 
\end{abstract}

\maketitle

\section{Introduction}	
In this paper, we study Open Problem~2.1 posed by Brezis in \cite{Brezis2023}, which concerns the uniqueness of solutions to the following Ginzburg--Landau equation:
\begin{equation}\label{eq:GL-system}
	\begin{cases}
		-\Delta u
		=
		\dfrac{1}{\varepsilon^2}u\bigl(1-\abs{u}^2\bigr)
		& \text{in } \Omega,\\[2mm]
		u(x)=x
		& \text{on } \partial\Omega,
	\end{cases}
\end{equation}
where $\Omega$ is the unit disc in $\mathbb R^2$, $\varepsilon>0$ is a given parameter, and the unknown $u$ maps $\Omega$ into $\mathbb R^2$. Originating from the theory of superconductivity, the Ginzburg--Landau equation has become a fundamental model in the study of vortices \cite{GinzburgLandau1950,deGennes1966}.

It is well known that \eqref{eq:GL-system} admits a radial solution of the form
\[	U_\varepsilon(x)
=
\frac{x}{\abs{x}}f_\varepsilon\bigl(\abs{x}\bigr),\]
where $f_\varepsilon\colon[0,1]\to\R$ satisfies
\[	\begin{cases}
	-f_\varepsilon''
	-\dfrac{1}{r}f_\varepsilon'
	+\dfrac{1}{r^2}f_\varepsilon
	=
	\dfrac{1}{\varepsilon^2}
	f_\varepsilon\bigl(1-f_\varepsilon^2\bigr)
	& \text{in }(0,1),\\[2mm]
	f_\varepsilon(0)=0,
	\qquad
	f_\varepsilon(1)=1.
\end{cases}\]
A natural question is whether the radial solution $U_\varepsilon$ is the unique solution of \eqref{eq:GL-system}, which is formulated as Open Problem~2.1 in \cite{Brezis2023}.


\begin{brezisproblem}
	Is the radial solution $U_\varepsilon$ the only solution of
	\eqref{eq:GL-system} for every $\varepsilon>0$?
\end{brezisproblem}
An affirmative answer to Brezis' Open Problem~2.1 would, in particular,
imply that every solution of the Ginzburg--Landau equation is radially
symmetric and coincides with the radial solution determined by the
boundary data. This is reminiscent of the classical Gidas--Ni--Nirenberg theorem for positive solutions of scalar elliptic equations
\cite{GidasNiNirenberg1979}. However, the vector-valued nature of the
Ginzburg--Landau equation prevents a direct use of the maximum-principle
and moving-plane methods underlying such scalar symmetry results.





Although several known results support an affirmative answer to Brezis' Open Problem~2.1 in \cite{Brezis2023}, uniqueness is not known for all $\varepsilon>0$. So far, uniqueness has been established in two parameter regimes. Let $\lambda_1$ denote the first Dirichlet eigenvalue of $-\Delta$ in $\Omega$. For $\varepsilon\geq1/\sqrt{\lambda_1}$, uniqueness follows from the strict convexity of the energy functional associated with equation \eqref{eq:GL-system}
\[
E_\varepsilon(u)
:=
\frac{1}{2}\int_{\Omega}\abs{\nabla u}^{2}
+
\frac{1}{4\varepsilon^{2}}
\int_{\Omega}\bigl(\abs{u}^{2}-1\bigr)^{2},
\]
defined on
\[
H_g^1(\Omega;\mathbb R^2)
:=
\left\{
u\in H^1(\Omega;\mathbb R^2):
u=g \text{ on }\partial\Omega
\right\},
\qquad
g(x)=x.
\]
Indeed, weak solutions of \eqref{eq:GL-system} are critical points of $E_\varepsilon$, and strict convexity therefore gives uniqueness; see Lemma~\ref{lem:convexity-threshold} in Section~3 for details. At the other end of the parameter range, Pacard and Rivi\`ere \cite{PacardRiviere2000} proved that $U_\varepsilon$ is also the unique solution when $\varepsilon>0$ is sufficiently small, namely, for $\varepsilon\in(0,\varepsilon_1)$ for some $\varepsilon_1>0$. Mironescu \cite{Mironescu1995} proved that, for every $\varepsilon>0$, the radial solution $U_\varepsilon$ is a local minimizer of $E_\varepsilon$ and that $D^2E_\varepsilon(U_\varepsilon)$ is positive definite. This local stability result, however, does not imply uniqueness among all solutions, since it does not exclude other critical points away from $U_\varepsilon$.



Other uniqueness results are known under additional assumptions. Comte and Mironescu in \cite{ComteMironescu1999} obtained uniqueness results under additional assumptions, with uniqueness proved for sufficiently small $\varepsilon$ under a non-vanishing assumption on the solution and for every $\varepsilon>0$ under suitable conditions on the boundary data. However, these assumptions do not cover the degree-one boundary condition
\[
g(x)=x\quad\text{on }\partial\Omega
\]
in problem \eqref{eq:GL-system}, and hence do not settle Brezis' Open Problem~2.1. Recently, Chen, Liu, Wei and Yang proved that $U_\varepsilon$ is a global minimizer of $E_\varepsilon$, thereby resolving Brezis' Open Problem~2.2; see \cite{ChenLiuWeiYang2026}. This result, however, does not imply uniqueness among all solutions of \eqref{eq:GL-system}, since a critical point need not be a global minimizer.




After decreasing $\varepsilon_1$ if necessary, we may assume that $\varepsilon_1<1/\sqrt{\lambda_1}$. Consequently, the remaining difficulty for Brezis' Open Problem~2.1 lies in the intermediate range
\[
\varepsilon\in\left(\varepsilon_1,\frac{1}{\sqrt{\lambda_1}}\right),
\]
where neither the small-$\varepsilon$ result nor the strict convexity argument applies.



Our main result provides a partial answer to Brezis' Open Problem~2.1 by extending the known uniqueness range below the threshold $1/\sqrt{\lambda_1}$. More precisely, we prove that uniqueness still holds for $\varepsilon<1/\sqrt{\lambda_1}$ when $\varepsilon$ is sufficiently close to $1/\sqrt{\lambda_1}$. Thus, uniqueness persists below the convexity threshold, into a parameter regime where the strict convexity argument is no longer available. Combined with the known uniqueness for $\varepsilon\geq1/\sqrt{\lambda_1}$, this yields a constant $\delta>0$ such that \eqref{eq:GL-system} admits a unique weak solution for every
\[
\varepsilon\in\left(\frac{1}{\sqrt{\lambda_1}}-\delta,\infty\right).
\]
In fact, this result follows from a more general uniqueness theorem for the Ginzburg--Landau equation with nontrivial boundary conditions, which includes Brezis' Open Problem~2.1 as a special case. Consider the following Ginzburg--Landau equation:
\begin{equation}\label{eq:general-GL-system}
	\begin{cases}
		-\Delta u
		=
		\dfrac{1}{\varepsilon^2}
		u\bigl(1-\abs{u}^2\bigr)
		& \text{in } \Omega,\\[2mm]
		u=g
		& \text{on } \partial\Omega,
	\end{cases}
\end{equation}
where $2\leq N\leq4$, $\Omega\subset\R^N$ is a bounded connected
$C^{1,1}$ domain, $\varepsilon>0$, $m\geq1$, 
$u\colon\Omega\to\R^m$  and $g$ satisfies
\[
	g\in W^{2-1/p,p}(\partial\Omega, \R^m),
	\qquad
	g\not\equiv0,
\]
for some fixed $p>N$. Unlike the uniqueness results mentioned above that require a non-vanishing solution or additional boundary assumptions, our theorem does not impose any non-vanishing condition on $u$; it only requires the boundary data to be nontrivial, namely, $g\not\equiv0$. In particular, the argument does not rely on radial symmetry or on the specific degree-one boundary condition in Brezis' problem 2.1. For the general problem \eqref{eq:general-GL-system}, we use the same notation
$\lambda_1=\lambda_1(\Omega)$ for the first Dirichlet eigenvalue
of $-\Delta$ in $\Omega$ and set
$\varepsilon_*:=1/\sqrt{\lambda_1}$.

We now state our main results.

\begin{theorem}\label{thm:general-uniqueness}
	There exists $\delta\in(0,\varepsilon_*)$ such that, for every
	\[
		\varepsilon\in
		\left(\varepsilon_*-\delta,\infty\right),
	\]
problem \eqref{eq:general-GL-system} has a unique weak solution.
\end{theorem}



The threshold $\varepsilon_*=\lambda_1^{-1/2}$ plays a distinguished role in the proof of Theorem~\ref{thm:general-uniqueness}. At $\varepsilon=\varepsilon_*$, the Poincar\'e inequality makes the quadratic form associated with the linearized operator nonnegative. If $h$ belongs to the kernel of the linearized operator, equality in the Poincar\'e inequality forces $h$ to have the form $h=\phi_1\xi$, where $\phi_1>0$ is a first Dirichlet eigenfunction and $\xi\in\mathbb R^m$. Since $g\not\equiv0$, the solution $u_*$ at $\varepsilon_*$ is nontrivial, and the additional potential term in the quadratic form then forces $\xi=0$. Hence the linearized operator has trivial kernel. To obtain the isomorphism property required by the implicit function theorem, we next prove coercivity of the associated bilinear form; the resulting weak solvability, together with elliptic regularity, yields surjectivity and hence invertibility of the linearized operator. The implicit function theorem then gives local uniqueness near $(\varepsilon_*,u_*)$.




Local uniqueness alone does not exclude solutions lying far from this implicit-function-theorem branch. To rule them out, we establish a compactness result for arbitrary solutions whose parameters approach $\varepsilon_*$: every sequence of such solutions has a subsequence converging to $u_*$ in $W^{2,p}$. Hence all solutions for parameters sufficiently close to $\varepsilon_*$ must lie in the local uniqueness neighborhood. Combining this fact with strict convexity for $\varepsilon\geq\varepsilon_*$ yields Theorem~\ref{thm:general-uniqueness}.



Theorem~\ref{thm:general-uniqueness}, applied with $N=m=2$,
$\Omega$ the unit disc, and $g(x)=x$, gives the following partial
answer to Brezis' Open Problem~2.1.

\begin{corollary}\label{thm:local-uniqueness}
	There exists $\delta>0$ such that, for every
	\[
	\varepsilon\in
	\left(
	\varepsilon_*-\delta,
	\infty
	\right),
	\]
	the radial solution $U_\varepsilon$ is the unique solution of
	\eqref{eq:GL-system}.
\end{corollary}


The paper is organized as follows. Section~2 is devoted to the variational
framework and the existence of weak solutions for problem \eqref{eq:general-GL-system}. In Section~3, we prove the uniqueness of weak
solutions. The case $\varepsilon\geq\varepsilon_*$ follows from the
strict convexity of the energy functional, while uniqueness for $\varepsilon<\varepsilon_*$ sufficiently close to $\varepsilon_*$ is obtained by a local analysis around
the solution at $\varepsilon=\varepsilon_*$, combined with compactness and
the implicit function theorem.

\section{Variational setting and existence result}

In this section, we set up the variational framework for \eqref{eq:general-GL-system} and prove the existence of weak solutions in the general setting of Theorem~\ref{thm:general-uniqueness}. Although the existence of solutions to \eqref{eq:GL-system} in Brezis' Open Problem~2.1 is already known, the corresponding existence result for the more general problem \eqref{eq:general-GL-system} needs to be established. To this end, we exploit the variational structure of the Ginzburg--Landau equation and obtain weak solutions as the critical points of the associated energy functional.

Due to the nonhomogeneous boundary condition, the admissible functions
form an affine space rather than a linear space. A natural way to handle
this is to introduce a harmonic extension of the boundary data $g$. By
\cite[Theorem~1.5.1.2]{Grisvard1985}, there exists
$G\in W^{2,p}(\Omega;\R^m)$ such that $G=g$ on $\partial\Omega$.


Applying \cite[Theorem~9.15]{GilbargTrudinger2001}, there exists a unique
$w\in W^{2,p}(\Omega;\R^m)\cap W_0^{1,p}(\Omega;\R^m)$ satisfying
\[
	\begin{cases}
		-\Delta w=\Delta G & \text{in }\Omega,\\
		w=0 & \text{on }\partial\Omega.
	\end{cases}
\]
Then $\widetilde g:=G+w\in W^{2,p}(\Omega;\R^m)$ is the harmonic
extension of $g$, namely, it satisfies
\[
	\begin{cases}
		-\Delta\widetilde g=0 & \text{in }\Omega,\\
		\widetilde g=g & \text{on }\partial\Omega.
	\end{cases}
\]
Therefore, the admissible space can be written as
\[
	H_g^1(\Omega;\R^m)
	:=
	\widetilde g+H_0^1(\Omega;\R^m).
\]

The energy functional associated with the Ginzburg--Landau equation is given by
\[E_\varepsilon(u)
:=
\frac{1}{2}\int_\Omega\abs{\nabla u}^2
+
\frac{1}{4\varepsilon^2}
\int_\Omega\bigl(\abs{u}^2-1\bigr)^2, \qquad
u\in H_g^1(\Omega;\R^m).\]
Since $2\leq N\leq4$, the Sobolev
embedding
$
	H^1(\Omega)\hookrightarrow L^4(\Omega)
$
(see, e.g., \cite{Brezis2011}) ensures that $E_\varepsilon$
is well defined on $H_g^1(\Omega;\R^m)$ and is of class $C^2$.

A direct computation gives, for
$u\in H_g^1(\Omega;\R^m)$,
\begin{equation}\label{eq:first-variation}
	DE_\varepsilon(u)[\varphi]
	=
	\int_\Omega
	\nabla u:\nabla\varphi
	-
	\frac{1}{\varepsilon^2}
	\int_\Omega
	\bigl(1-\abs{u}^2\bigr)u\cdot\varphi
	\qquad
	\forall\,\varphi\in H_0^1(\Omega;\R^m),
\end{equation}
and
\begin{equation}\label{eq:second-variation}
	\begin{aligned}
		D^2E_\varepsilon(u)[\varphi,\psi]
		={}&
		\int_\Omega
		\nabla\varphi:\nabla\psi
		-
		\frac{1}{\varepsilon^2}
		\int_\Omega
		\bigl(1-\abs{u}^2\bigr)\varphi\cdot\psi\\
		&+
		\frac{2}{\varepsilon^2}
		\int_\Omega
		(u\cdot\varphi)(u\cdot\psi)
	\end{aligned}
	\qquad
	\forall\,\varphi,\psi\in H_0^1(\Omega;\R^m).
\end{equation}

\begin{definition}
	A function $u\in H_g^1(\Omega;\R^m)$ is called a weak solution of
	\eqref{eq:general-GL-system} if
	\[
	\int_\Omega
	\nabla u:\nabla\varphi
	=
	\frac{1}{\varepsilon^2}
	\int_\Omega
	\bigl(1-\abs{u}^2\bigr)u\cdot\varphi,
	\qquad
	\forall\,\varphi\in H_0^1(\Omega;\R^m).
	\]
\end{definition}
For $\varepsilon>0$, denote by $\mathcal{S}_\varepsilon$ the set
of all weak solutions of \eqref{eq:general-GL-system} with parameter
$\varepsilon$:
\begin{equation}\label{eq:solution-set}
	\mathcal{S}_\varepsilon
	:=
	\left\{
	u\in H_g^1(\Omega;\R^m):
	u \text{ is a weak solution of }
	\eqref{eq:general-GL-system}
	\right\}.
\end{equation}

We now establish the existence of weak solutions for \eqref{eq:general-GL-system}.

\begin{lemma}\label{lem:existence}
	For every $\varepsilon>0$, problem
	\eqref{eq:general-GL-system} admits a weak solution
	$u\in H_g^1(\Omega;\R^m)$.
\end{lemma}

\begin{proof}
	Let $\varepsilon>0$ be fixed. Since $E_\varepsilon\geq0$ on
	$H_g^1(\Omega;\R^m)$, we have
	\[
		0
		\leq
		c_\varepsilon
		:=
		\inf_{u\in H_g^1(\Omega;\R^m)}E_\varepsilon(u)
		\leq
		E_\varepsilon(\widetilde g)
		<
		\infty.
	\]
	Let $(u_n)\subset H_g^1(\Omega;\R^m)$ be such that
	$E_\varepsilon(u_n)\to c_\varepsilon$. For $n$ sufficiently large,
	\begin{equation}\label{eq:existence-gradient-bound}
		\frac{1}{2}
		\norm{\nabla u_n}_{L^2(\Omega)}^2
		\leq
		E_\varepsilon(u_n)
		\leq
		c_\varepsilon+1.
	\end{equation}
	Set $v_n:=u_n-\widetilde g\in H_0^1(\Omega;\R^m)$. By the
	Poincar\'e inequality,
	\begin{equation}\label{eq:existence-poincare-bound}
		\norm{v_n}_{L^2(\Omega)}
		\leq
		C\norm{\nabla v_n}_{L^2(\Omega)}
		\leq
		C\left(
		\norm{\nabla u_n}_{L^2(\Omega)}
		+
		\norm{\nabla\widetilde g}_{L^2(\Omega)}
		\right).
	\end{equation}
	Together with \eqref{eq:existence-gradient-bound} and
	\eqref{eq:existence-poincare-bound}, $(u_n)$ is bounded in
	$H^1(\Omega;\R^m)$.

	Since $H^1(\Omega;\R^m)$ is reflexive, after passing to a
	subsequence, there exists $u\in H^1(\Omega;\R^m)$ such that
	\[
		u_n\rightharpoonup u
		\qquad\text{in }H^1(\Omega;\R^m)
		\quad\text{as }n\to\infty.
	\]
Since $u_n-\widetilde g\in H_0^1(\Omega;\R^m)$ and
$u_n-\widetilde g\rightharpoonup u-\widetilde g$ in
$H^1(\Omega;\R^m)$, the weak closedness of $H_0^1(\Omega;\R^m)$ implies
that $u-\widetilde g\in H_0^1(\Omega;\R^m)$. Hence
$u\in H_g^1(\Omega;\R^m)$. Since $\Omega$ is bounded, the embedding
$H^1(\Omega)\hookrightarrow L^2(\Omega)$ is compact. Thus, after
passing to a further subsequence,
	\[
		u_n\to u
		\qquad\text{in }L^2(\Omega;\R^m)
		\quad\text{as }n\to\infty,
	\]
	and
	\[
		u_n(x)\to u(x)
		\qquad\text{for a.e. }x\in\Omega
		\quad\text{as }n\to\infty.
	\]
	By weak lower semicontinuity,
	\[
		\int_\Omega\abs{\nabla u}^2
		\leq
		\liminf_{n\to\infty}
		\int_\Omega\abs{\nabla u_n}^2.
	\]
	Moreover, by Fatou's lemma,
	\[
		\int_\Omega
		\bigl(\abs{u}^2-1\bigr)^2
		\leq
		\liminf_{n\to\infty}
		\int_\Omega
		\bigl(\abs{u_n}^2-1\bigr)^2.
	\]
	Therefore,
	\[
		E_\varepsilon(u)
		\leq
		\liminf_{n\to\infty}E_\varepsilon(u_n)
		=
		c_\varepsilon.
	\]
	Since $u\in H_g^1(\Omega;\R^m)$, the definition of $c_\varepsilon$
	also gives $c_\varepsilon\leq E_\varepsilon(u)$. Hence
	$E_\varepsilon(u)=c_\varepsilon$, and $u$ is a minimizer of
	$E_\varepsilon$ on $H_g^1(\Omega;\R^m)$. Thus
	$DE_\varepsilon(u)[\varphi]=0$ for every
	$\varphi\in H_0^1(\Omega;\R^m)$, and
	\eqref{eq:first-variation} shows that $u$ is a weak solution of
	\eqref{eq:general-GL-system}. This proves the lemma.
\end{proof}

\section{Uniqueness results}

In this section, we turn to the uniqueness problem. As mentioned in the
introduction, strict convexity of the associated energy functional yields
uniqueness for \eqref{eq:general-GL-system} when
$\varepsilon\geq\varepsilon_*$, where
$\varepsilon_*=\lambda_1^{-1/2}$. Thus it remains to consider the case
$\varepsilon<\varepsilon_*$, where strict convexity is no longer
available. We study the behavior of solutions near $\varepsilon_*$ and
prove that there exists $\delta>0$ such that the weak solution is unique
for
\[
\varepsilon\in(\varepsilon_*-\delta,\infty).
\]
Throughout this section, we set
$\varepsilon_*=\lambda_1^{-1/2}$ and denote by $u_*$ the unique weak
solution of \eqref{eq:general-GL-system} corresponding to
$\varepsilon=\varepsilon_*$. We use $\mathcal{S}_\varepsilon$ as defined
in \eqref{eq:solution-set}.


\subsection{Uniqueness for $\varepsilon\geq\varepsilon_*$}


\begin{lemma}\label{lem:convexity-threshold}
	For every $\varepsilon\geq\varepsilon_*$, the functional
	$E_\varepsilon$ is strictly convex on
	$H_g^1(\Omega;\R^m)$. Consequently,
	problem \eqref{eq:general-GL-system} has a unique weak solution.
\end{lemma}

\begin{proof}
	Let $u\in H_g^1(\Omega;\R^m)$ and
	$h\in H_0^1(\Omega;\R^m)\setminus\{0\}$. By
	\eqref{eq:second-variation},
	\[
		\begin{aligned}
		D^2E_\varepsilon(u)[h,h]
		={}&
		\int_\Omega\abs{\nabla h}^2
		-
		\frac{1}{\varepsilon^2}
		\int_\Omega\abs{h}^2\\
		&+
		\frac{1}{\varepsilon^2}
		\int_\Omega\abs{u}^2\abs{h}^2
		+
		\frac{2}{\varepsilon^2}
		\int_\Omega(u\cdot h)^2.
		\end{aligned}
	\]
	Since $\varepsilon\geq\varepsilon_*$, the Poincar\'e inequality gives
	\begin{equation}\label{eq:convexity-poincare}
		\int_\Omega\abs{\nabla h}^2
		-
		\frac{1}{\varepsilon^2}
		\int_\Omega\abs{h}^2
		\geq
		\left(
		\lambda_1-\frac{1}{\varepsilon^2}
		\right)
		\int_\Omega\abs{h}^2
		\geq0.
	\end{equation}
	Suppose that $D^2E_\varepsilon(u)[h,h]=0$. By
	\eqref{eq:second-variation} and \eqref{eq:convexity-poincare},
	\[
		\int_\Omega\abs{\nabla h}^2
		-
		\frac{1}{\varepsilon^2}
		\int_\Omega\abs{h}^2
		=
		0,
		\qquad
		\int_\Omega\abs{u}^2\abs{h}^2
		=
		0.
	\]
	Since $h\neq0$, \eqref{eq:convexity-poincare} gives
	$\varepsilon=\varepsilon_*$ and
	\[
		\int_\Omega\abs{\nabla h}^2
		=
		\lambda_1\int_\Omega\abs{h}^2.
	\]
	Since $\Omega$ is connected, the first Dirichlet eigenvalue
	$\lambda_1$ is simple. Moreover, for every
	$w=(w_1,\ldots,w_m)\in H_0^1(\Omega;\R^m)$,
	\[
		\int_\Omega\abs{\nabla w}^2
		-
		\lambda_1\int_\Omega\abs{w}^2
		=
		\sum_{\alpha=1}^m
		\left(
		\int_\Omega\abs{\nabla w_\alpha}^2
		-
		\lambda_1\int_\Omega\abs{w_\alpha}^2
		\right).
	\]
	Each term on the right-hand side is nonnegative by the Poincar\'e
	inequality. Therefore,
	\begin{equation}\label{eq:first-eigenspace-characterization}
		\int_\Omega\abs{\nabla w}^2
		=
		\lambda_1\int_\Omega\abs{w}^2
		\quad\Longrightarrow\quad
		w=\phi_1\xi
		\quad\text{for some }\xi\in\R^m,
	\end{equation}
	where $\phi_1$ is the positive Dirichlet eigenfunction corresponding
	to $\lambda_1$. Applying
	\eqref{eq:first-eigenspace-characterization} to $h$, we have
	$h=\phi_1\xi$ for some $\xi\in\R^m\setminus\{0\}$. It follows that
	\[
		0
		=
		\int_\Omega\abs{u}^2\abs{h}^2
		=
		\abs{\xi}^2
		\int_\Omega\abs{u}^2\phi_1^2.
	\]
	Since $\phi_1>0$ in $\Omega$, we obtain $u=0$ a.e. in $\Omega$,
	which contradicts $u\in H_g^1(\Omega;\R^m)$ and $g\not\equiv0$.
	Therefore $D^2E_\varepsilon(u)[h,h]>0$ for every
	$u\in H_g^1(\Omega;\R^m)$ and every
	$h\in H_0^1(\Omega;\R^m)\setminus\{0\}$. Let
	$u_0,u_1\in H_g^1(\Omega;\R^m)$ with $u_0\neq u_1$, and define
	$q\colon[0,1]\to\R$ by
	$q(t):=E_\varepsilon(u_0+t(u_1-u_0))$. Then
	\[
		q''(t)
		=
		D^2E_\varepsilon
		\bigl(u_0+t(u_1-u_0)\bigr)
		[u_1-u_0,u_1-u_0]
		>
		0
	\]
	for every $t\in[0,1]$. Hence $q$ is strictly convex, and therefore
	$E_\varepsilon$ is strictly convex on $H_g^1(\Omega;\R^m)$.
	By Lemma~\ref{lem:existence}, a weak solution exists for every
	$\varepsilon>0$. Moreover, weak solutions are precisely the critical
	points of $E_\varepsilon$ on $H_g^1(\Omega;\R^m)$, and a
	differentiable strictly convex functional has at most one critical
	point. Consequently, \eqref{eq:general-GL-system} has a unique weak
	solution for every $\varepsilon\geq\varepsilon_*$.

\end{proof}

\subsection{Proof of Theorem~\ref{thm:general-uniqueness}}


In this subsection, we prove Theorem~\ref{thm:general-uniqueness}. When
$\varepsilon<\varepsilon_*$, the strict convexity argument used above is
no longer available. Therefore, we cannot obtain uniqueness directly from
the convexity of $E_\varepsilon$. We first study the limiting behavior of
solutions as $\varepsilon\to\varepsilon_*$. The following lemma shows that
any sequence of solutions whose parameters converge to $\varepsilon_*$
admits a subsequence converging to $u_*$ in $W^{2,p}$. This compactness
property is crucial in the proof of Theorem~\ref{thm:general-uniqueness}.


\begin{lemma}\label{lem:compactness}
	Let $\varepsilon_n\to\varepsilon_*$ and
	$u_n\in\mathcal{S}_{\varepsilon_n}$. Then, after passing to a
	subsequence,
	\[
		u_n\to u_*
		\qquad\text{in }W^{2,p}(\Omega;\R^m)
		\quad\text{as }n\to\infty.
	\]
\end{lemma}

\begin{proof}
	Since $p>N$, the Sobolev embedding
	$W^{2-1/p,p}(\partial\Omega;\R^m)
	\hookrightarrow L^\infty(\partial\Omega;\R^m)$ gives
	$g\in L^\infty(\partial\Omega;\R^m)$. Set
	$M:=\max\{1,\norm{g}_{L^\infty(\partial\Omega)}\}$. We first establish the estimate
	\begin{equation}\label{eq:bound-un}
		\abs{u_n}\leq M
		\qquad\text{a.e. in }\Omega.
	\end{equation}
	Define
	\[
	T(z):=
	\begin{cases}
		0, & \abs{z}\leq M,\\[1mm]
		\left(1-\dfrac{M^2}{\abs{z}^2}\right)z,
		& \abs{z}>M.
	\end{cases}
	\]
	Since $T$ is globally Lipschitz and $T(g)=0$ on
	$\partial\Omega$, we have $T(u_n)\in H_0^1(\Omega;\R^m)$.
	Let $A_n:=\{x\in\Omega:\abs{u_n(x)}>M\}$. Using $T(u_n)$
	as a test function, we obtain
	\begin{align*}
		\int_{A_n}
		\nabla u_n:\nabla T(u_n)
		&=
		\frac{1}{\varepsilon_n^2}
		\int_{A_n}
		\bigl(1-\abs{u_n}^2\bigr)u_n\cdot T(u_n)\\
		&=
		\frac{1}{\varepsilon_n^2}
		\int_{A_n}
		\bigl(1-\abs{u_n}^2\bigr)
		\bigl(\abs{u_n}^2-M^2\bigr).
	\end{align*}
	On the other hand, for a.e. $x\in A_n$,
\[
\nabla u_n:\nabla T(u_n)
=
\left(1-\frac{M^2}{\abs{u_n}^2}\right)\abs{\nabla u_n}^2
+
\frac{2M^2}{\abs{u_n}^4}
\sum_{j=1}^N
\bigl(u_n\cdot\partial_j u_n\bigr)^2
\geq0.
\]
\begin{equation}\label{eq:truncation-estimate}
	0
	\leq
	\int_{A_n}
	\nabla u_n:\nabla T(u_n)
	=
	\frac{1}{\varepsilon_n^2}
	\int_{A_n}
	\bigl(1-\abs{u_n}^2\bigr)
	\bigl(\abs{u_n}^2-M^2\bigr)
	\leq0.
\end{equation}
	Since $M\geq1$,
	$(1-\abs{u_n}^2)(\abs{u_n}^2-M^2)<0$ on $A_n$.
	Hence \eqref{eq:truncation-estimate} gives $\abs{A_n}=0$, and
	\eqref{eq:bound-un} follows.
	
	Set
	\[
	f_n:=
	\frac{1}{\varepsilon_n^2}
	\bigl(1-\abs{u_n}^2\bigr)u_n.
	\]
	Since $\varepsilon_n\to\varepsilon_*$ as $n\to\infty$ and
	$\varepsilon_*>0$, it follows from \eqref{eq:bound-un} that
	there exists a constant $C_1>0$, independent of $n$, such that
	\[
	\norm{f_n}_{L^\infty(\Omega)}\leq C_1.
	\]
	Since $\Omega$ is bounded, we deduce that
	\[
	\norm{f_n}_{L^p(\Omega)}
	\leq
	\abs{\Omega}^{1/p}\norm{f_n}_{L^\infty(\Omega)}
	\leq C_2,
	\]
	where $C_2:=\abs{\Omega}^{1/p}C_1$.
	
	Let $v_n:=u_n-\widetilde g$.
	Then $v_n\in H_0^1(\Omega;\R^m)$ and
	\[
	\begin{cases}
		-\Delta v_n=f_n & \text{in }\Omega,\\
		v_n=0 & \text{on }\partial\Omega.
	\end{cases}
	\]
	By \cite[Theorem~9.15]{GilbargTrudinger2001} and uniqueness of the
	weak solution,
	\[
	v_n\in W^{2,p}(\Omega;\R^m)\cap
	W_0^{1,p}(\Omega;\R^m).
	\]
	In addition, by \cite[Lemma~9.17]{GilbargTrudinger2001}, there exists
	$C=C(N,p,\Omega)>0$, independent of $n$, such that
	\[
	\norm{v_n}_{W^{2,p}(\Omega)}
	\leq
	C\norm{f_n}_{L^p(\Omega)}
	\leq
	CC_2.
	\]
	
	Consequently,
	\[
	\norm{u_n}_{W^{2,p}(\Omega)}
	\leq
	CC_2+\norm{\widetilde g}_{W^{2,p}(\Omega)}
	=:C_3.
	\]
	Hence $(u_n)$ is bounded in $W^{2,p}(\Omega;\R^m)$.
	

Since $p>N$, for every $0<\alpha<1-\frac{N}{p}$ the embedding
	\[
	W^{2,p}(\Omega)\hookrightarrow
	C^{1,\alpha}(\overline{\Omega})
	\]
	is compact (see \cite[Theorem~7.26(ii)]{GilbargTrudinger2001}).
	Thus, after passing to a subsequence, there exists
	$u\in W^{2,p}(\Omega;\R^m)$ such that
	\begin{equation}\label{eq:compact-convergence}
	u_n\rightharpoonup u
	\quad\text{in }W^{2,p}(\Omega;\R^m),
	\qquad
	u_n\to u
	\quad\text{in }C^{1,\alpha}(\overline{\Omega};\R^m),
	\end{equation}
	as $n\to\infty$. By \eqref{eq:compact-convergence} and
	$u_n=g$ on $\partial\Omega$, we have $u=g$ on $\partial\Omega$.
	
	Let $G(z):=(1-\abs{z}^2)z$. By
	\eqref{eq:compact-convergence}, there exists $R>0$ such that
	\[
	\norm{u_n}_{L^\infty(\Omega)}\leq R
	\qquad\text{and}\qquad
	\norm{u}_{L^\infty(\Omega)}\leq R
	\]
	for every $n$. Since $G$ is Lipschitz continuous on
	$\overline{B_R(0)}$, there exists a constant $C_1>0$ such that
	\[
	\abs{G(z)-G(w)}
	\leq C_1\abs{z-w}
	\qquad
	\text{for all }z,w\in\overline{B_R(0)}.
	\]
	Consequently,
	\[
	\begin{aligned}
		\norm{G(u_n)-G(u)}_{L^p(\Omega)}
		&\leq C_1\norm{u_n-u}_{L^p(\Omega)}\\
		&\leq C_2\norm{u_n-u}_{L^\infty(\Omega)}
		\longrightarrow0,
	\end{aligned}
	\]
	where $C_2:=C_1\abs{\Omega}^{1/p}$.
	Therefore,
	\[
	\frac{1}{\varepsilon_n^2}G(u_n)
	\to
	\frac{1}{\varepsilon_*^2}G(u)
	\qquad\text{in }L^p(\Omega;\R^m)
	\quad\text{as }n\to\infty.
	\]
For every $\varphi\in H_0^1(\Omega;\R^m)$, we have
\[
\int_\Omega
\nabla u_n:\nabla\varphi
=
\frac{1}{\varepsilon_n^2}
\int_\Omega G(u_n)\cdot\varphi.
\]
Passing to the limit as $n\to\infty$, we obtain
\[
\int_\Omega
\nabla u:\nabla\varphi
=
\frac{1}{\varepsilon_*^2}
\int_\Omega G(u)\cdot\varphi,
\qquad
\forall\,\varphi\in H_0^1(\Omega;\R^m).
\]
	Since $u=g$ on $\partial\Omega$, we have
	$u\in\mathcal{S}_{\varepsilon_*}$. By
	Lemma~\ref{lem:convexity-threshold}, $u=u_*$.
	
	Set $w_n:=u_n-u_*$. Then
	$w_n\in W^{2,p}(\Omega;\R^m)\cap
	W_0^{1,p}(\Omega;\R^m)$ and
	\[
	-\Delta w_n
	=
	\frac{1}{\varepsilon_n^2}G(u_n)
	-
	\frac{1}{\varepsilon_*^2}G(u_*).
	\]
	Moreover,
	\[
	\frac{1}{\varepsilon_n^2}G(u_n)
	-
	\frac{1}{\varepsilon_*^2}G(u_*)
	\to0
	\qquad\text{in }L^p(\Omega;\R^m)
	\quad\text{as }n\to\infty.
	\]
	By \cite[Lemma~9.17]{GilbargTrudinger2001},
	\[
	\begin{aligned}
		\norm{u_n-u_*}_{W^{2,p}(\Omega)}
		&=
		\norm{w_n}_{W^{2,p}(\Omega)}\\
		&\leq
		C
		\norm{
			\frac{1}{\varepsilon_n^2}G(u_n)
			-
			\frac{1}{\varepsilon_*^2}G(u_*)
		}_{L^p(\Omega)}
		\longrightarrow0
	\end{aligned}
	\]
	as $n\to\infty$.
	This proves the lemma.
\end{proof}

Having established the compactness of solutions as
$\varepsilon\to\varepsilon_*$, we next analyze the local structure near
$u_*$. To this end, we formulate \eqref{eq:general-GL-system} as a
nonlinear operator equation and study its linearization at
$(\varepsilon_*,u_*)$.

Let
\[
X:=W^{2,p}(\Omega;\R^m)\cap W_0^{1,p}(\Omega;\R^m),
\qquad
Y:=L^p(\Omega;\R^m),
\]
where $X$ is equipped with the norm
\[
\norm{u}_{W^{2,p}(\Omega)}
=
\left(
\sum_{\abs{\alpha}\leq2}
\norm{D^\alpha u}_{L^p(\Omega)}^p
\right)^{1/p}.
\]
Let
$v_*:=u_*-\widetilde g$ and define
\[
\mathcal{F}\colon(0,\infty)\times X\to Y
\]
by
\[
\mathcal{F}(\varepsilon,v)
:=
-\Delta v
-\frac{1}{\varepsilon^2}
\bigl(1-\abs{\widetilde g+v}^2\bigr)
(\widetilde g+v).
\]
Since $p>N$, $W^{2,p}(\Omega)$ embeds continuously into
$C^{1,\alpha}(\overline{\Omega})$, and hence into $L^\infty(\Omega)$,
for every $0<\alpha<1-\frac{N}{p}$; see
\cite[Theorem~7.26(ii)]{GilbargTrudinger2001}. 
Since
$z\mapsto(1-\abs{z}^2)z$ is smooth, it follows that
$\mathcal{F}\in C^1((0,\infty)\times X,Y)$. The truncation argument in
the proof of Lemma~\ref{lem:compactness} gives
$u\in L^\infty(\Omega;\R^m)$ for every
$u\in\mathcal{S}_\varepsilon$. Since
$u-\widetilde g\in H_0^1(\Omega;\R^m)$ and
\[
-\Delta(u-\widetilde g)
=
\frac{1}{\varepsilon^2}(1-\abs{u}^2)u
\in L^p(\Omega;\R^m),
\]
elliptic regularity gives
$u\in W^{2,p}(\Omega;\R^m)$. Hence,
$\mathcal{F}(\varepsilon,v)=0$ if and only if
$\widetilde g+v\in\mathcal{S}_\varepsilon$. In particular,
$\mathcal{F}(\varepsilon_*,v_*)=0$.



A direct computation gives
\[
D_v\mathcal{F}(\varepsilon_*,v_*)h
=
-\Delta h
-\lambda_1\bigl(1-\abs{u_*}^2\bigr)h
+
2\lambda_1(u_*\cdot h)u_*.
\]
We denote this operator by
\[
L_*:=D_v\mathcal{F}(\varepsilon_*,v_*)\colon X\to Y.
\]

The following nondegeneracy property is crucial for applying the implicit function theorem and obtaining local uniqueness near $u_*$.

\begin{lemma}\label{lem:nondegeneracy}
	The operator $L_*\colon X\to Y$ is an isomorphism.
\end{lemma}
\begin{proof}
The operator $L_*\colon X\to Y$ is linear. By
	Lemma~\ref{lem:compactness}, $u_*\in W^{2,p}(\Omega;\R^m)$, and hence
	$u_*\in L^\infty(\Omega;\R^m)$. Moreover, there exists
	$C>0$ such that
	\[
	\norm{L_*h}_{L^p(\Omega)}
	\leq
	\norm{\Delta h}_{L^p(\Omega)}
	+
	\lambda_1
	\left(
	1+3\norm{u_*}_{L^\infty(\Omega)}^2
	\right)
	\norm{h}_{L^p(\Omega)}
	\leq
	C\norm{h}_{W^{2,p}(\Omega)}.
	\]
	Thus $L_*\colon X\to Y$ is bounded. It remains to prove that $L_*$
	is bijective; boundedness of the inverse then follows from the bounded
	inverse theorem.

We first prove that $L_*$ is injective. Let $h\in X$ satisfy
	$L_*h=0$. Testing this equation with $h$ and integrating by parts, we obtain
	\begin{equation}\label{eq:kernel-identity}
	0
	=
	\int_\Omega\abs{\nabla h}^2
	-\lambda_1\int_\Omega\abs{h}^2
	+\lambda_1\int_\Omega\abs{u_*}^2\abs{h}^2
	+2\lambda_1\int_\Omega(u_*\cdot h)^2.
	\end{equation}
	By the Poincar\'e inequality,
	\[
	\int_\Omega\abs{\nabla h}^2
	-\lambda_1\int_\Omega\abs{h}^2
	\geq0.
	\]
	Combining this with \eqref{eq:kernel-identity}, we obtain
	\[
		\int_\Omega\abs{\nabla h}^2
		=
		\lambda_1\int_\Omega\abs{h}^2,
		\qquad
		\int_\Omega\abs{u_*}^2\abs{h}^2=0.
	\]
	By \eqref{eq:first-eigenspace-characterization},
	$h=\phi_1\xi$ for some $\xi\in\R^m$. Therefore,
	\begin{equation}\label{eq:kernel-eigenfunction-identity}
	0
	=
	\int_\Omega\abs{u_*}^2\abs{h}^2
	=
	\abs{\xi}^2
	\int_\Omega\abs{u_*}^2\phi_1^2.
	\end{equation}
	By the boundary condition $u_*=g$ on $\partial\Omega$ and
	$g\not\equiv0$, we have $u_*\not\equiv0$. Since
	$\phi_1>0$ in $\Omega$,
	\eqref{eq:kernel-eigenfunction-identity} implies $\xi=0$.
	Thus $h=0$, and then $L_*$ is injective.
	
	We next prove that $L_*$ is surjective. On
	$H_0^1(\Omega;\R^m)$, consider the bilinear form
	\[
	a(h,k)
	:=
	\int_\Omega
	\nabla h:\nabla k
	-\lambda_1\int_\Omega
	\bigl(1-\abs{u_*}^2\bigr)h\cdot k
	+2\lambda_1\int_\Omega
	(u_*\cdot h)(u_*\cdot k).
	\]
	By the Poincaré inequality and
	$u_*\in L^\infty(\Omega;\R^m)$, there exists $C>0$ such that
	\[
	\abs{a(h,k)}
	\leq
	C
	\norm{\nabla h}_{L^2(\Omega)}
	\norm{\nabla k}_{L^2(\Omega)}
	\]
	for every $h,k\in H_0^1(\Omega;\R^m)$, and thus the form $a$ is
	bounded. By the Poincar\'e inequality,
	\[
		\begin{aligned}
		a(h,h)
		={}&
		\int_\Omega\abs{\nabla h}^2
		-\lambda_1\int_\Omega\abs{h}^2
		+\lambda_1\int_\Omega\abs{u_*}^2\abs{h}^2\\
		&+
		2\lambda_1\int_\Omega(u_*\cdot h)^2
		\geq0
		\end{aligned}\]
	for every $h\in H_0^1(\Omega;\R^m)$.
	We argue by contradiction. Suppose that $a$ is not coercive. Then for
	every $n\geq1$ there would exist
	$\widetilde{h}_n\in H_0^1(\Omega;\R^m)\setminus\{0\}$ such that
	\[
	a(\widetilde{h}_n,\widetilde{h}_n)
	<
	\frac{1}{n}
	\norm{\nabla\widetilde{h}_n}_{L^2(\Omega)}^2.
	\]
	Define
	\[
	h_n
	:=
	\frac{\widetilde{h}_n}
	{\norm{\nabla\widetilde{h}_n}_{L^2(\Omega)}}.
	\]
	Then
	\[
	\norm{\nabla h_n}_{L^2(\Omega)}=1
	\qquad\text{and}\qquad
	0\leq a(h_n,h_n)<\frac{1}{n}.
	\]
	Hence $a(h_n,h_n)\to0$ as $n\to\infty$.
	By the definition of $a$,
	\begin{align*}
		a(h_n,h_n)
		={}&
		\int_\Omega\abs{\nabla h_n}^2
		-\lambda_1\int_\Omega\abs{h_n}^2\\
		&+
		\lambda_1\int_\Omega
		\abs{u_*}^2\abs{h_n}^2
		+
		2\lambda_1\int_\Omega
		(u_*\cdot h_n)^2,
	\end{align*}
	while the Poincar\'e inequality gives
	\[
	\int_\Omega\abs{\nabla h_n}^2
	-\lambda_1\int_\Omega\abs{h_n}^2
	\geq0.
	\]
	Together with $a(h_n,h_n)\to0$, this yields
	\[
	\int_\Omega\abs{\nabla h_n}^2
	-\lambda_1\int_\Omega\abs{h_n}^2
	\to0
	\qquad\text{as }n\to\infty
	\]
	and
	\[
	\int_\Omega\abs{u_*}^2\abs{h_n}^2
	\to0
	\qquad\text{as }n\to\infty.
	\]
	Since $\norm{\nabla h_n}_{L^2(\Omega)}=1$, the sequence
	$(h_n)$ is bounded in $H_0^1(\Omega;\R^m)$. By the reflexivity of
	$H_0^1(\Omega;\R^m)$ and the compact embedding
	$H_0^1(\Omega)\hookrightarrow L^2(\Omega)$, after passing to a
	subsequence, there exists $h\in H_0^1(\Omega;\R^m)$ such that
	\begin{equation}\label{eq:coercivity-convergence}
	h_n\rightharpoonup h
	\quad\text{in }H_0^1(\Omega;\R^m),
	\qquad
	h_n\to h
	\quad\text{in }L^2(\Omega;\R^m)
	\quad\text{as }n\to\infty.
	\end{equation}
	Since $\norm{\nabla h_n}_{L^2(\Omega)}=1$ and
	$\int_\Omega\abs{\nabla h_n}^2
	-\lambda_1\int_\Omega\abs{h_n}^2\to0$,
	\eqref{eq:coercivity-convergence} gives
	$\lambda_1\int_\Omega\abs{h}^2=1$.
	By weak lower semicontinuity,
	\[
	\int_\Omega\abs{\nabla h}^2
	\leq
	\liminf_{n\to\infty}
	\int_\Omega\abs{\nabla h_n}^2
	=1.
	\]
	On the other hand, the Poincaré inequality gives
	\[
	\int_\Omega\abs{\nabla h}^2
	\geq
	\lambda_1\int_\Omega\abs{h}^2
	=1.
	\]
	By \eqref{eq:first-eigenspace-characterization},
	$h=\phi_1\xi$ for some $\xi\in\R^m\setminus\{0\}$. Since
	$h_n\to h$ in $L^2(\Omega;\R^m)$ and
	$u_*\in L^\infty(\Omega;\R^m)$, we also have
	\[
	\int_\Omega\abs{u_*}^2\abs{h_n}^2
	\to
	\int_\Omega\abs{u_*}^2\abs{h}^2.
	\]
	Consequently,
	\[
	\int_\Omega\abs{u_*}^2\abs{h}^2=0,
	\]
	which contradicts $\phi_1>0$ in $\Omega$ and
	$u_*\not\equiv0$. Hence there exists $c>0$ such that
	\[
	a(h,h)
	\geq
	c\norm{\nabla h}_{L^2(\Omega)}^2
	\qquad
	\text{for every }h\in H_0^1(\Omega;\R^m).
	\]
	
	Let $f\in Y$. Since $p>N\geq2$ and $\Omega$ is bounded, we have
	$f\in L^2(\Omega;\R^m)$, and the map
	\[
	k\longmapsto\int_\Omega f\cdot k
	\]
	is a bounded linear functional on $H_0^1(\Omega;\R^m)$. By
	\cite[Theorem~5.8]{GilbargTrudinger2001}, there exists a unique
	$h\in H_0^1(\Omega;\R^m)$ such that
	\[
	a(h,k)=\int_\Omega f\cdot k
	\qquad
	\text{for every }k\in H_0^1(\Omega;\R^m).
	\]
	Thus $h$ is a weak solution of
	\[
	-\Delta h
	=
	f
	+
	\lambda_1\bigl(1-\abs{u_*}^2\bigr)h
	-
	2\lambda_1(u_*\cdot h)u_*.
	\]
	Since $f\in L^2(\Omega;\R^m)$,
	$u_*\in L^\infty(\Omega;\R^m)$, and
	$h\in L^2(\Omega;\R^m)$, the right-hand side belongs to
	$L^2(\Omega;\R^m)$.
	By \cite[Theorem~9.15]{GilbargTrudinger2001} and uniqueness of the
	weak solution, we have $h\in W^{2,2}(\Omega;\R^m)$.

	If $2\leq N<4$, choose $r\in(N/2,2)$. Since $\Omega$ is bounded, the embedding
	\[
		W^{2,2}(\Omega)\hookrightarrow W^{2,r}(\Omega)
	\]
	is continuous. By \cite[Theorem~7.26(ii)]{GilbargTrudinger2001},
	\[
		W^{2,r}(\Omega)
		\hookrightarrow
		C^{0,\alpha}(\overline{\Omega})
		\hookrightarrow
		L^\infty(\Omega),
		\qquad
		\alpha=2-\frac{N}{r}>0.
	\]
	If $N=4$, set $r:=2p/(p+2)$. Since $p>4$, we have
	$1<r<2$ and $2r<4$. Moreover,
	\[
		W^{2,2}(\Omega)\hookrightarrow W^{2,r}(\Omega)
	\]
	and \cite[Theorem~7.26(i)]{GilbargTrudinger2001} gives
	\[
		W^{2,r}(\Omega)
		\hookrightarrow
		L^{4r/(4-2r)}(\Omega)
		=
		L^p(\Omega).
	\]
	Thus, in every case, $h\in L^p(\Omega;\R^m)$. Together with
	$f\in L^p(\Omega;\R^m)$ and
	$u_*\in L^\infty(\Omega;\R^m)$, this gives
	$f+\lambda_1(1-\abs{u_*}^2)h
	-2\lambda_1(u_*\cdot h)u_*\in L^p(\Omega;\R^m)$. By
	\cite[Theorem~9.15]{GilbargTrudinger2001} and uniqueness of the weak
	solution,
	\[
		h\in W^{2,p}(\Omega;\R^m)\cap W_0^{1,p}(\Omega;\R^m)=X.
	\]
	Therefore $L_*h=f$, and $L_*$ is surjective. It follows from
	\cite[Corollary~2.12(c)]{Rudin1991} that $L_*^{-1}\colon Y\to X$ is
	bounded. This proves the lemma.
\end{proof}

We are now in a position to prove
Theorem~\ref{thm:general-uniqueness}.

\begin{proof}[Proof of Theorem~\ref{thm:general-uniqueness}]
	By Lemma~\ref{lem:nondegeneracy} and the implicit function theorem
	\cite[Theorem~15.1 and Corollary~15.1]{Deimling1985}, there exist
	neighborhoods $I\subset(0,\infty)$ of $\varepsilon_*$ and
	$V\subset X$ of $v_*$, and a unique $C^1$ map
	$T\colon I\to V$ such that $T(\varepsilon_*)=v_*$ and
	\[
		\mathcal{F}(\varepsilon,v)=0
		\quad\Longleftrightarrow\quad
		v=T(\varepsilon)
	\]
	for every $(\varepsilon,v)\in I\times V$.


We claim that there exists $\delta_0>0$ such that
	\eqref{eq:general-GL-system} has at most one weak solution for every
	$\varepsilon\in(\varepsilon_*-\delta_0,\varepsilon_*)$.
	Arguing by contradiction, we find a sequence
	$\varepsilon_n\in(0,\varepsilon_*)$ with
	$\varepsilon_n\to\varepsilon_*$ and two distinct solutions
	$u_n,\widehat{u}_n\in\mathcal{S}_{\varepsilon_n}$.



	By Lemma~\ref{lem:compactness}, after passing to a subsequence,
	\[
		u_n\to u_*,
		\qquad
		\widehat{u}_n\to u_*
		\qquad\text{in }W^{2,p}(\Omega;\R^m).
	\]


	Set $v_n:=u_n-\widetilde g$ and
	$\widehat{v}_n:=\widehat{u}_n-\widetilde g$. Then
	\[
		v_n\to v_*,
		\qquad
		\widehat{v}_n\to v_*
		\qquad\text{in }X.
	\]
	Thus, for $n$ sufficiently large,
	$(\varepsilon_n,v_n),(\varepsilon_n,\widehat{v}_n)\in I\times V$.
	Since
	$\mathcal{F}(\varepsilon_n,v_n)=
	\mathcal{F}(\varepsilon_n,\widehat{v}_n)=0$, the local uniqueness provided
	by the implicit function theorem yields
	\[
		v_n=T(\varepsilon_n)=\widehat{v}_n.
	\]
	Hence $u_n=\widehat{u}_n$, which is a contradiction. This proves
	the claim.


By Lemma~\ref{lem:existence}, the claim gives a unique weak solution for
	every $\varepsilon\in(\varepsilon_*-\delta_0,\varepsilon_*)$, while
	Lemma~\ref{lem:convexity-threshold} gives uniqueness for every
	$\varepsilon\geq\varepsilon_*$. After decreasing $\delta_0$ if
	necessary so that $\delta_0<\varepsilon_*$ and setting
	$\delta=\delta_0$, we obtain the conclusion of
	Theorem~\ref{thm:general-uniqueness}.

\end{proof}

\begin{proof}[Proof of Corollary~\ref{thm:local-uniqueness}]
	Apply Theorem~\ref{thm:general-uniqueness} with $N=m=2$,
	$\Omega$ the unit disc, and $g(x)=x$.
	Since $U_\varepsilon$ is a weak solution of \eqref{eq:GL-system}, it is
	the unique weak solution for every
	$\varepsilon\in(1/\sqrt{\lambda_1}-\delta,\infty)$, for some
	$\delta>0$.
\end{proof}

\subsection*{Conflict of interest}

The authors declare no conflict of interest.

\subsection*{Ethics approval}
 Not applicable.

\subsection*{Data Availability Statements}
Data sharing not applicable to this article as no datasets were generated or analysed during the current study.

\subsection*{Acknowledgements}
 C. Ji is partially supported by National Natural Science Foundation of China (No. 12571117).

\begin{thebibliography}{99}

\bibitem{Brezis2011} 
H.~Brezis, 
\textit{Functional Analysis, Sobolev Spaces and Partial Differential Equations}, 
Springer, New York, 2011. 

\bibitem{Brezis2023} 
H.~Brezis, 
\textit{Some of my favorite open problems}, 
Rend.~Lincei Mat.~Appl. \textbf{34} (2023), 307--335.	 

\bibitem{ChenLiuWeiYang2026} 
H.-G.~Chen, Y.~Liu, J.~Wei and W.~Yang, 
\textit{On Brezis' Open Problem 2.2}, 
arXiv:2608.15957 [math.AP], 2026. 

\bibitem{ComteMironescu1999} 
M.~Comte and P.~Mironescu, 
\textit{Minimizing properties of arbitrary solutions to the Ginzburg--Landau equation}, 
Proc.~Roy.~Soc.~Edinburgh Sect.~A \textbf{129} (1999), no.~6, 1157--1169.

\bibitem{deGennes1966} 
P.-G.~de Gennes, 
\textit{Superconductivity of Metals and Alloys}, 
W.~A.~Benjamin, New York, 1966. 

\bibitem{Deimling1985} 
K.~Deimling, 
\textit{Nonlinear functional analysis}, 
Springer-Verlag, Berlin, 1985.	 

\bibitem{GidasNiNirenberg1979} 
B.~Gidas, W.~M.~Ni and L.~Nirenberg, 
\textit{Symmetry and related properties via the maximum principle}, 
Commun.~Math.~Phys.~\textbf{68} (1979), 209--243. 

\bibitem{GilbargTrudinger2001} 
D.~Gilbarg and N.~S.~Trudinger, 
\textit{Elliptic partial differential equations of second order}, 
Classics in Mathematics, Springer-Verlag, Berlin, 2001. 

\bibitem{GinzburgLandau1950} 
V.~L.~Ginzburg and L.~D.~Landau, 
\textit{On the theory of superconductivity}, 
Zh.~Eksp.~Teor.~Fiz. \textbf{20} (1950), 1064--1082. 

\bibitem{Grisvard1985} 
P.~Grisvard, 
\textit{Elliptic problems in nonsmooth domains}, 
Monographs and Studies in Mathematics, vol.~24, 
Pitman (Advanced Publishing Program), Boston, MA, 1985. 

\bibitem{Mironescu1995}
P.~Mironescu,
\textit{On the stability of radial solutions of the Ginzburg--Landau equation},
J.~Funct.~Anal. \textbf{130} (1995), no.~2, 334--344.


\bibitem{PacardRiviere2000} 
F.~Pacard and T.~Rivi\`ere, 
\textit{Linear and nonlinear aspects of vortices: 
	the Ginzburg--Landau model}, 
Progress in Nonlinear Differential Equations and Their Applications, 
vol.~39, Birkh\"auser Boston, Boston, MA, 2000. 

\bibitem{Rudin1991} 
W.~Rudin, 
\textit{Functional analysis}, 
2nd ed., International Series in Pure and Applied Mathematics, 
McGraw-Hill, Inc., New York, 1991.

\end{thebibliography}	
\end{document}
